% GENERATED FILE. Edit canonical lessons, then run npm run generate:books.
% canonical-source-hash: fnv1a64:10ea9c4888b9c330
% canonical-lessons: ZH-C213-qianzheng, ZH-C213-daoyou, ZH-C213-youke, ZH-C213-menpiao, ZH-C213-yilupingan

\chapter{A Trip Abroad}
\label{ch:zh-a-trip-abroad}
\hlchaptermodality{213}

\begin{chapteropening}
\textbf{By the end of this chapter:} \emph{I can talk about visas, tour guides, tourists and tickets, and wish someone a safe journey.}

Complete the last lesson of chapter 213: I can talk about visas, tour guides, tourists and tickets, and wish someone a safe journey.
\end{chapteropening}

\section[qiānzhèng]{qiānzhèng — visa}

\label{lesson:ZH-C213-qianzheng}



\begin{quote}
Before the new one: say the Mandarin for to make an appointment, then the Mandarin for to get a haircut.
\end{quote}

\subsection*{You'll want to know: qiānzhèng}
\emph{qiānzhèng} — \textquotedblleft{}visa\textquotedblright{}.

\begin{quote}
\emph{wǒ hái méiyǒu qiānzhèng} — I haven't got a visa yet.
\end{quote}

Its characters come later; for now it is a word for the ear.

\begin{cousinweb}
\emph{qiānzhèng} is built from \emph{qiān} 'to sign' and \emph{zhèng} 'proof'. The second is inherited from Middle Chinese \emph{tsyingH} (Guangyun, in Baxter's transcription), glossed 'to verify'; no dictionary we use records an older form of the first, so its history is not settled.
\end{cousinweb}

\begin{grammarlens}[title={What you've built}]
One more word for talking about plans.
\end{grammarlens}

\subsection*{Your turn}
\begin{itemize}
\raggedright
  \item \emph{Say it:} \emph{qiānzhèng}
  \item \emph{Say it:} \emph{qiānzhèng}, once more
  \item \emph{Say it:} \emph{wǒ hái méiyǒu qiānzhèng} — I haven't got a visa yet.
  \item \emph{Recall:} say the Mandarin for to make an appointment, then the Mandarin for to get a haircut, then say \emph{qiānzhèng} again
\end{itemize}

\begin{tcolorbox}[breakable,colback=teal!4,colframe=teal!35!black,title={Before you move on}]
What does qiānzhèng mean? (\textquotedblleft{}Visa\textquotedblright{}.) Say it once more.
\end{tcolorbox}
\section[dǎoyóu]{dǎoyóu — tour guide}

\label{lesson:ZH-C213-daoyou}



\begin{quote}
Before the new one: say the Mandarin for to get a haircut, then the Mandarin for visa.
\end{quote}

\subsection*{You'll want to know: dǎoyóu}
\emph{dǎoyóu} — \textquotedblleft{}tour guide\textquotedblright{}.

\begin{quote}
\emph{wǒmen de dǎoyóu shì běndì rén} — Our tour guide is a local.
\end{quote}

Its characters come later; for now it is a word for the ear.

\begin{cousinweb}
\emph{dǎoyóu} is built from \emph{dǎo} 'to lead' and \emph{yóu} 'to travel about', the \emph{yóu} of \emph{lǚyóu}, both inherited from Middle Chinese, \emph{dawH} and \emph{yuw} (Guangyun, in Baxter's transcription), which glosses the first 'to lead'. The dictionary gives it a departing tone, so the rising tone of dǎo is not the regular outcome, and how it came is not settled.
\end{cousinweb}

\begin{grammarlens}[title={What you've built}]
One more word for talking about plans.
\end{grammarlens}

\subsection*{Your turn}
\begin{itemize}
\raggedright
  \item \emph{Say it:} \emph{dǎoyóu}
  \item \emph{Say it:} \emph{dǎoyóu}, once more
  \item \emph{Say it:} \emph{wǒmen de dǎoyóu shì běndì rén} — Our tour guide is a local.
  \item \emph{Recall:} say the Mandarin for to get a haircut, then the Mandarin for visa, then say \emph{dǎoyóu} again
\end{itemize}

\begin{tcolorbox}[breakable,colback=teal!4,colframe=teal!35!black,title={Before you move on}]
What does dǎoyóu mean? (\textquotedblleft{}Tour guide\textquotedblright{}.) Say it once more.
\end{tcolorbox}
\section[yóukè]{yóukè — tourist}

\label{lesson:ZH-C213-youke}



\begin{quote}
Before the new one: say the Mandarin for visa, then the Mandarin for tour guide.
\end{quote}

\subsection*{You'll want to know: yóukè}
\emph{yóukè} — \textquotedblleft{}tourist\textquotedblright{}.

\begin{quote}
\emph{zhèr de yóukè dōu shì wàiguó rén} — The tourists here are all from abroad.
\end{quote}

Its characters come later; for now it is a word for the ear.

\begin{cousinweb}
\emph{yóukè} is built from \emph{yóu} 'to travel about' and \emph{kè} 'a guest', both inherited from Middle Chinese, \emph{yuw} and \emph{khaek} (Guangyun, in Baxter's transcription).
\end{cousinweb}

\begin{grammarlens}[title={What you've built}]
One more word for talking about plans.
\end{grammarlens}

\subsection*{Your turn}
\begin{itemize}
\raggedright
  \item \emph{Say it:} \emph{yóukè}
  \item \emph{Say it:} \emph{yóukè}, once more
  \item \emph{Say it:} \emph{zhèr de yóukè dōu shì wàiguó rén} — The tourists here are all from abroad.
  \item \emph{Recall:} say the Mandarin for visa, then the Mandarin for tour guide, then say \emph{yóukè} again
\end{itemize}

\begin{tcolorbox}[breakable,colback=teal!4,colframe=teal!35!black,title={Before you move on}]
What does yóukè mean? (\textquotedblleft{}Tourist\textquotedblright{}.) Say it once more.
\end{tcolorbox}
\section[\texorpdfstring{\zh{门票} (ménpiào)}{门票 (ménpiào)}]{\zh{门票} (ménpiào) — admission ticket}

\label{lesson:ZH-C213-menpiao}



\begin{quote}
Before the new one: say the Mandarin for tour guide, then the Mandarin for tourist.
\end{quote}

\subsection*{You'll want to know: \zh{门票}}
\textbf{\zh{门票}} — \emph{ménpiào} — \textquotedblleft{}admission ticket\textquotedblright{}.

\begin{quote}
\textbf{\zh{门票多少钱}} \emph{ménpiào duōshao qián} — How much is it to get in?
\end{quote}

\begin{cousinweb}
\textbf{\zh{门票}} (\emph{ménpiào}) is built from \zh{门} 'gate' and \zh{票} 'a ticket'. The first is inherited from Middle Chinese \emph{mwon} (Guangyun, in Baxter's transcription); no dictionary we use records an older form of \zh{票}, so its history is not settled.
\end{cousinweb}

\begin{grammarlens}[title={What you've built}]
One more word for talking about plans.
\end{grammarlens}

\subsection*{Your turn}
\begin{itemize}
\raggedright
  \item \emph{Say it:} \emph{ménpiào}
  \item \emph{Say it:} \emph{ménpiào}, once more
  \item \emph{Say it:} \emph{ménpiào duōshao qián} — How much is it to get in?
  \item \emph{Recall:} say the Mandarin for tour guide, then the Mandarin for tourist, then say \emph{ménpiào} again
\end{itemize}

\begin{tcolorbox}[breakable,colback=teal!4,colframe=teal!35!black,title={Before you move on}]
What does \zh{门票} mean? (\textquotedblleft{}Admission ticket\textquotedblright{}.) Say it once more.
\end{tcolorbox}
\section[yílù píng'ān]{yílù píng'ān — have a safe journey}

\label{lesson:ZH-C213-yilupingan}



\begin{quote}
Before the new one: say the Mandarin for tourist, then the Mandarin for admission ticket.
\end{quote}

\subsection*{You'll want to know: yílù píng'ān}
\emph{yílù píng'ān} — \textquotedblleft{}have a safe journey\textquotedblright{}.

\begin{quote}
\emph{nǐ míngtiān huíguó yílù píng'ān} — You're going home tomorrow: have a safe journey!
\end{quote}

Its characters come later; for now it is a word for the ear.

\zh{一} changes its tone here: yí, before the falling tone of lù.

\begin{cousinweb}
\emph{yílù píng'ān} is built from \zh{一路} 'the whole way' and \emph{píng'ān} 'safe and sound', all four parts inherited from Middle Chinese, \emph{'jit}, \emph{luH}, \emph{bjaeng} and \emph{'an} (Guangyun, in Baxter's transcription).
\end{cousinweb}

\begin{grammarlens}[title={What you've built}]
One more word for talking about plans.
\end{grammarlens}

\subsection*{Your turn}
\begin{itemize}
\raggedright
  \item \emph{Say it:} \emph{yílù píng'ān}
  \item \emph{Say it:} \emph{yílù píng'ān}, once more
  \item \emph{Say it:} \emph{nǐ míngtiān huíguó yílù píng'ān} — You're going home tomorrow: have a safe journey!
  \item \emph{Recall:} say the Mandarin for tourist, then the Mandarin for admission ticket, then say \emph{yílù píng'ān} again
\end{itemize}

\begin{tcolorbox}[breakable,colback=teal!4,colframe=teal!35!black,title={Before you move on}]
What does yílù píng'ān mean? (\textquotedblleft{}Have a safe journey\textquotedblright{}.) Say it once more.
\end{tcolorbox}
